\title{Byte Latent Transformer: Patches Scale Better Than Tokens}

\begin{abstract}
We present the Byte Latent Transformer (\textsc{BLT}), a novel byte-level large language model architecture that, for the first time, achieves tokenization-based LLM performance at scale while offering notable gains in inference speed and robustness. \textsc{BLT} transforms bytes into adaptively determined patches, which become the fundamental computational units of the model. These patches are created by segmenting input according to the entropy of the subsequent byte, concentrating computational effort and model capacity where the data's complexity rises. We provide the inaugural flop-controlled scaling analysis of byte-level models, spanning up to 8B parameters and utilizing 4T training bytes. Our findings confirm that it is practical to scale models that learn directly from unprocessed bytes, dispensing with any fixed vocabulary. Training and inference are made more efficient by dynamically choosing larger patches in regions of predictable data, which also delivers qualitative gains in reasoning and enhanced generalization to rare cases. Overall, when controlling for inference cost, \textsc{BLT} demonstrates much improved scaling compared to models relying on tokenization, by simultaneously expanding both patch granularity and model size.
\end{abstract}

\section{Introduction}
\label{section:intro}

We present the Byte Latent Transformer~(\textbf{BLT}), a novel architecture that dispenses with tokenization altogether by directly learning from unprocessed byte sequences. This approach, for the first time, achieves parity in performance with models reliant on tokenization when scaled, while also offering marked gains in both efficiency and robustness (\S\ref{sec:robustness}).
In contrast, current large language models (llms) are trained in an almost entirely end-to-end fashion, with the notable exception of tokenization—a rule-based preprocessing operation that converts byte streams into a predetermined token set. This practice imposes structural biases in sequence compression, resulting in drawbacks like heightened sensitivity to domain or modality~\citep{dagangetting}, increased vulnerability to input perturbations~(\S\ref{sec:robustness}), limited encoding of orthographic information~\citep{edman2024cute}, and persistent inequities for multilingual data~\citep{liang2023xlm,petrov2024language,limisiewicz2024myte}.

In the past, tokenization has been crucial, as training large language models directly on bytes proves extremely expensive at scale because of the extended sequence lengths involved~\citep{xue2022byt5}. To address this, earlier studies have utilized more efficient self-attention mechanisms~\citep{el2020characterbert, clark2022canine} or employed architectures that bypass attention altogether~\citep{wang2024mambabyte}~(\S\ref{section:related_work}). Nonetheless, these strategies are primarily effective for training \textit{small models}. For large-scale models, most of the computational overhead in a Transformer arises from the substantial feed-forward network layers applied at each byte, rather than from the attention component itself.

In order to optimize compute distribution, we introduce an adaptive, trainable approach to aggregating bytes into \textit{patches}~(\S\ref{section:patching}), alongside a novel model structure that integrates both byte-level and patch-level data. Distinct from traditional tokenization, {\textsc BLT} does not depend on a predetermined patch vocabulary. Instead, any combination of bytes can be transformed into latent patch embeddings through efficient, parameterized encoder and decoder components. Our findings indicate that this strategy enables a \textit{more} effective use of computational resources compared to tokenization-centric models.

Most tokenization-based language models dedicate identical computational resources to every token. This approach compromises performance for the sake of efficiency, as the heuristics used for compression do not consistently align with the intricacy of the prediction tasks. A foundational principle of our model is that computational effort should be adaptively distributed, directing more resources to harder tasks as necessary. For instance, employing a sizable transformer to generate the ending of a typical word is superfluous, given that such predictions are generally straightforward, low-entropy tasks—unlike the selection of the opening word in a new sentence, which is more challenging. This principle informs the design of {\textsc BLT}'s architecture~(\S\ref{section:architecture}), which incorporates three transformer modules: a pair of compact byte-level \textit{local models} and a sizable, global \textit{latent transformer}~(\autoref{fig:arch_high_level}).

In order to decide how bytes should be aggregated into patches—and hence how computational resources are dynamically apportioned—{\textsc BLT} utilizes segmentation based on the entropy associated with next-byte predictions. This process leads to groupings of bytes that share similar informational content within their contexts.

We introduce the first systematic scaling analysis of byte-level models under a fixed flop budget, extending up to 8B parameters and 4T bytes of training data, and demonstrate the feasibility of large-scale, end-to-end byte-based training without reliance on fixed-vocabulary tokenization. In terms of training flop efficiency, {\textsc BLT} achieves comparable results to Llama 3, yet delivers up to 50\% reduction in inference computational requirements~(\S\ref{section:scaling}), as measured by Floating Point OPerations (flops).\footnote{We estimate model computation by enumerating Floating Point OPerations (flops).} 

Additionally, we find that direct modeling over raw bytes confers notable advantages, particularly for data distribution tails. {\textsc BLT} consistently surpasses tokenizer-dependent baselines in robustness to input noise and exhibits superior fine-grained character-level understanding—substantiated through tasks targeting orthography, phonological features, and scenarios of low-resource machine translation~(\S\ref{sec:robustness}). 

Moreover, {\textsc BLT} enables concurrent expansion of both model size and patch size without exceeding a given inference flop allowance. Employing larger patch sizes generally reduces computational cost, thus making it possible to invest more resources in scaling up the global latent transformer, as its application frequency diminishes. Our inference-flop controlled scaling trials~(\autoref{fig:fixed_inference_scaling}) further reveal that byte-level architectures display significantly improved scaling behavior in comparison to tokenization-based approaches.

In conclusion, our work offers several key contributions: 1) We propose {\textsc BLT}, a byte latent llm framework that adaptively distributes computational resources to enhance flop efficiency, 2) We demonstrate that {\textsc BLT} can reach flop-matched training performance comparable to Llama 3 at the 8B scale, while also providing the flexibility to exchange slight reductions in evaluation scores for substantial increases in flop efficiency—up to 50\%, 3) Our {\textsc BLT} models introduce a novel approach to scaling llms, enabling model size expansion without exceeding a predetermined inference budget, and 4) We provide evidence that {\textsc BLT} models exhibit superior resilience to input noise and recognize sub-word characteristics in input data, a capability lacking in token-based llms. The source code for training and inference with {\textsc BLT} is publicly available at~\url{https://github.com/facebookresearch/blt}.

\section{Patching: From Individual Bytes to Groups of Bytes}
\label{section:patching}

Dividing a byte sequence into \textit{patches} enables {\textsc BLT} to flexibly allocate computational resources depending on the input’s characteristics. As illustrated in Figure~\ref{fig:patching_types}, there are multiple strategies for segmenting bytes into patches. More formally, a patching function $f_p$ takes a sequence of bytes $\pmb{x} = \{x_i | i=1, \ldots, n\}$ of length $n$ and divides it into a series of $m < n$ patches $\pmb{p} = \{p_j | j=1, \ldots, m\}$ by associating each $x_i$ with a value from \{0,1\}, where 1 denotes the onset of a new patch. The predominant contributor to computational expense in both token-based and patch-based architectures is the number of execution steps required by the core Transformer, which in {\textsc BLT} equates to the number of patches prescribed by the chosen patching function. As a result, the average number of bytes in each patch—or \textit{patch size}—serves as the principal determinant of data processing costs for any given patching function during model training and inference~(\S\ref{section:flops}).

We proceed by presenting three distinct patching functions: fixed-size byte patching~(\S\ref{section:static-patch}), patching based on whitespace delimiters~(\S\ref{section:white-patch}), and adaptive patching that utilizes entropy estimates from a lightweight byte language model~(\S\ref{section:dyn-patch}). We conclude by examining the method of incremental patching, as well as clarifying the distinctions between tokenization and patching~(\S\ref{section:bpe-patch}).

\subsection{Strided Patching Every K Bytes}
\label{section:static-patch}

A common approach for organizing bytes is to divide them into equally sized patches of length $k$, following the strategy utilized by MegaByte~\citep{yu2023megabyte}. This fixed-length segmentation is straightforward to implement during both training and inference, offers a direct method for adjusting the average patch length, and allows for clear management of the computational workload. Despite these advantages, there are notable limitations to this method. Most importantly, computational resources are not adaptively distributed according to content needs: this can result in inefficiency, such as allocating transformer step $j$ to predict inconsequential regions (e.g., code whitespace) or failing to provide adequate computational capacity for segments densely packed with important information, like mathematical expressions. Additionally, fixed-length patching can result in inconsistent and non-semantic splitting of analogous byte sequences, exemplified by instances where identical words are segmented differently.

\subsection{Space Patching}
\label{section:white-patch}

\citet{slagle2024spacebyte} introduces an intuitive yet powerful enhancement to strided patching by generating new patches at each space-like byte\footnote{
    Space-like bytes are classified as any byte that does not represent a Latin character, digit, or utf-8 continuation byte. Furthermore, it is required that every patch includes at least one byte that is not space-like. }, capitalizing on their role as boundaries of linguistic units in numerous languages. Through Space patching, a latent transformer operation (i.e., increased computational effort) is applied to model every individual word. This method guarantees consistent patching for words across sequences and ensures computational resources are focused on making challenging predictions, which frequently occur just after spaces. As an illustration, predicting the initial byte in the completion to the prompt “Who composed the Magic Flute?\uline{\hspace{.6em}}” is significantly more demanding than predicting subsequent bytes after ``M,'' as the first letter narrows down the plausible continuations, rendering the remainder of ``Mozart'' much more straightforward to infer. Nonetheless, space patching exhibits limitations when faced with all languages and application domains, and a notable drawback is its inability to flexibly adjust patch size. We proceed by presenting a novel patching technique, drawing upon the observation that the initial bytes of words tend to be the hardest to predict, while simultaneously offering an inherent mechanism for modulating patch sizes.

\subsection{Entropy Patching: Using Next-Byte Entropies from a Small Byte LM}
\label{section:dyn-patch}

Instead of utilizing a rule-based method like whitespace to determine boundaries, we adopt a data-driven strategy for detecting next-byte predictions with elevated uncertainty. Specifically, we propose \textit{entropy patching}, which leverages entropy calculations to establish patch boundaries.

We train a small byte-level auto-regressive language model on the training data for {\textsc BLT} and compute next byte entropies under the LM distribution $p_{e}$ over the byte vocabulary $\mathcal{V}$:

\begin{equation}
H(x_i) = \sum_{v \in \mathcal{V}} p_{e}(x_i=v|\pmb{x}_{<i}) \log p_{e}(x_i=v|\pmb{x}_{<i})
\end{equation}

We explore two techniques for detecting patch boundaries utilizing entropies $H(x_i)$. The initial method selects points that exceed a predetermined global entropy threshold, as depicted in~\autoref{fig:patching}. The alternative method determines points where the entropy is substantially higher compared to the preceding value. This latter strategy can also be viewed as pinpointing locations where the generally monotonically decreasing entropy within a patch is disrupted.

\begin{align*}
    \text{Global Constraint}\hspace{1cm}H(x_t)& > \theta_g\\
    \text{Approx. Monotonic Constraint}\hspace{1cm}H(x_t)& - H(x_{t-1}) > \theta_r
\end{align*}

Patch boundaries are determined through an efficient preprocessing process carried out at the data loading stage. This approach contrasts with~\citet{nawrot-etal-2023-efficient}, who utilize a classifier trained to detect patch boundaries based on entropy metrics. In our study~(\S\ref{section:experiments}), we evaluate and contrast these two techniques for separating bytes of low and high entropy.

\subsection{The Byte-Pair Encoding (BPE) Tokenizer and Incremental Patching}
\label{section:incremental-patching}
\label{section:bpe-patch}

A large number of contemporary llms, such as our reference model Llama 3, rely on a subword tokenizer such as bpe~\citep{gage1994new, sennrich-etal-2016-neural}. In this work, we define ``tokens'' as groups of bytes selected from a \textit{finite} pre-established vocabulary set before training, distinguishing them from ``patches,'' which are adaptively formed sequences lacking a predetermined vocabulary. One key distinction is that, when using tokens, the model does not interact directly with the raw byte-level features.

An essential advancement offered by {\textsc BLT} relative to tokenization-based architectures lies in its reconfiguration of the balance between vocabulary size and computational requirements. Conventional llms experience a situation in which enlarging the vocabulary leads to tokens that, on average, represent larger segments of text, resulting in fewer processing steps for each sequence; however, this also causes the output layer of the model to expand due to the increased dimensionality tied to the enlarged vocabulary. Consequently, tokenization-based techniques face stringent constraints regarding their ability to markedly alter token granularity or inference demands. As an illustration, Llama 3 pushes the average token size from 3.7 to 4.4 bytes but does so at the expense of quadrupling the embedding table's size relative to Llama 2.

During generation, {\textsc BLT} must determine at each position in the byte sequence whether it has reached a patch boundary, as this impacts when additional computation via the Latent Transformer is triggered. This determination must rely solely on the already-generated portion of the sequence, without access to subsequent bytes. Consequently, the process of establishing patch boundaries must be feasible without foreknowledge of future bytes. In precise terms, a patching method $f_p$ adheres to the principle of incremental patching if and only if:
$$f_p(\pmb{x}_{<i}) = f_p(\pmb{x})_{<i}$$
bpe fails to be an incremental patching scheme, since an identical prefix may receive different tokenizations based on the sequence's continuation, violating the condition described above\footnote{Inserting a dedicated boundary token to demarcate patches can render bpe incremental; however, this comes at the cost of lengthening the byte sequence.}.

\section{{\textsc BLT} Architecture}
\label{section:architecture}
The {\textsc BLT} framework consists of three main components: a global autoregressive language model trained on patch-level representations, and two lightweight local modules responsible for encoding byte sequences into patches and reconstructing byte sequences from patch representations, respectively~(\autoref{fig:arch_high_level}).

\subsection{Latent Global Transformer Model}

\textit{The Latent Global Transformer} is an autoregressive transformer architecture, denoted as $\mathcal{G}$, composed of $l_{\mathcal{G}}$ layers. It processes a sequence of latent input patch embeddings, $p_j$, and transforms them into corresponding output patch embeddings, $o_j$. In this work, subscripts $j$ and $i$ are consistently used to refer to patches and bytes, respectively. For its attention mechanism, the global model implements a block-causal attention mask~\citep{dubey2024llama}, thereby limiting attention to encompass only the current patch and all prior patches within the document. This component dominates the computational cost during both pre-training and inference. Consequently, selectively triggering the global model enables dynamic adjustment of compute usage depending on the complexity of particular regions of the input and output.

\subsection{Local Encoder}
\label{section:local}

\textit{The Local Encoder Model}, indicated by $\mathcal{E}$, consists of a compact transformer architecture with $l_{\mathcal{E}} << l_{\mathcal{G}}$ layers. Its principal function is to rapidly convert an input byte sequence $b_i$ into rich patch-level representations, $p_j$. Notably, this architecture distinguishes itself from typical transformers by integrating a cross-attention layer after every transformer layer, serving to aggregate the individual byte representations into patch-level features~(\autoref{fig:crossattn}).

Initially, input bytes $b_i$ are projected into embeddings $x_i$ through a learnable matrix of shape $\mathbb{R}^{256 \times h_{\mathcal{E}}}$. Optionally, these embeddings can be enhanced with supplemental features provided by hash-embeddings~(\S\ref{sec:hash-emb}). 

The architecture then applies a sequence of alternating transformer and cross-attention blocks~(\S\ref{sec:enc-cross-attn}) to these embeddings, ultimately generating the patch representations $p_i$ that are subsequently input to the global transformer, $\mathcal{G}$. 

Within the transformer layers, a \textit{local block causal} attention scheme is adopted. Each byte attends to a fixed window of $w_{\mathcal{E}}$ preceding bytes, where this window can generally extend past evolving patch divisions, but does not cross the boundaries of the overall document. The subsequent sections offer a more detailed exposition of both the embedding approach and the structure of the cross-attention block.

\subsubsection{Encoder Hash n-gram Embeddings}
\label{sec:ngram-emb}
\label{sec:hash-emb}
An essential aspect of building strong, informative representations at each position $i$ involves leveraging contextual cues from preceding bytes. {\textsc BLT} addresses this by representing each byte $b_i$ not only on its own but also within the context of a byte n-gram. At every step $i$, we begin by forming byte-grams

\begin{align}
    g_{i,n}=\{b_{i-n + 1},\ldots, b_i\}
\end{align}

for every byte location $i$ and for $n$ ranging from three to eight.\footnote{
When $i<n$, byte-grams of size $n$ or greater are excluded. }

Next, we propose the use of hash $n$-gram embeddings, which involve mapping each possible byte $n$-gram, using a hash function, to a slot in a dedicated embedding matrix $E_{n}^{hash}$ with predetermined capacity, independently for each $n$ in the set $\{3,4,5,6,7,8\}$~\citep{bai2010learning}. These retrieved embeddings are summed with the original byte embedding, after which the aggregate representation is normalized and input to the local encoder network. The enriched byte embedding is given by:

\begin{align}
e_i &= x_i + \sum_{n = 3,...,8}E_{n}^{hash}(\text{Hash}(g_{i,n})) \\
\text{where, Hash}(g_{i,n}) &= \text{RollPolyHash}(g_{i,n}) \% |E_{n}^{hash}|
\end{align}

To adjust for variability in $n$-gram lengths, we divide $e_i$ by the total number of $n$-gram sizes plus one, employing RollPolyHash as specified in Appendix~\ref{appendix:rollpolyhash}. Section~\ref{section:ablations} examines how different choices of $n$ and the size of the embedding table affect the relationship between $n$-gram hash embeddings and flop-controlled scaling law patterns. Beyond utilizing hashed $n$-gram embeddings, we also assessed embeddings derived from $n$-gram frequency, with additional methodological specifics presented in Appendix~\ref{appendix:ngrams}.

\subsubsection{Encoder Multi-Headed Cross-Attention}
\label{sec:enc-cross-attn}
Our approach largely mirrors the cross-attention mechanism found in the input module of the Perceiver architecture~\citep{jaegle2021perceiver}. The key distinction lies in the way latent variables are defined: rather than employing a predetermined, fixed set of latent vectors, our method utilizes patch-specific latent representations that are dynamically determined for each patch~(\autoref{fig:crossattn}). Each latent representation only attends to the bytes within its associated patch. For each patch $p_j$, the module generates a query vector by pooling the representations of bytes included in $p_j$ and subsequently applying a linear projection using $\mathcal{E}_{C} \in \mathbb{R}^{h_{\mathcal{E}} \times (h_{\mathcal{E}} \times U_{\mathcal{E}})}$, where $U_{\mathcal{E}}$ signifies the total number of cross-attention heads in the encoder. Specifically, letting $f_{\text{bytes}}(p_j)$ represent the ordered list of bytes belonging to patch $p_j$, we compute

\begin{align}
P_{0,j} &= \mathcal{E}_{C}(f_\text{bytes}((p_j)), f~ \text{is a pooling function} \\
P_l &= P_{l-1} + W_o\left(\text{softmax}\left(\frac{QK^T}{\sqrt{d_k}}\right)V\right) \\
\text{where } Q_j &= W_q(P_{l-1,j}), K_i = W_k(h_{l-1,i}), V_i = W_v(h_{l-1,i}) \\
h_l &= \text{Encoder-Transformer-Layer}_l(h_{l-1})
\end{align}

Here, $P \in \mathbb{R}^{n_p \times h_{\mathcal{G}}}$ denotes the set of $n_p$ patch representations that are input to the global model, with each patch representation initially obtained by aggregating the byte embeddings $e_i$ associated with its respective patch $p_j$. The matrices $W_q$, $W_k$, $W_v$, and $W_o$ serve as linear projections for the queries, keys, values, and output, respectively; specifically, the keys and values are computed as projections of the byte-level representations $h_i$ from the previous network layer (or $e_i$ in the case of the first layer). To ensure each query $Q_j$ attends only to its own patch, a patch-based masking strategy is adopted: each $Q_j$ interacts exclusively with keys and values derived from bytes within patch $j$. As the model employs multi-headed attention on $Q, K$, and $V$—and since patch representations usually have a higher dimensionality ($h_{\mathcal{G}}$) than $h_{\mathcal{E}}$—the representation $P_l$ is structured as several attention heads of dimension $h_{\mathcal{E}}$ during cross-attention, which are subsequently concatenated to match the full dimension $h_{\mathcal{G}}$. Moreover, the model applies pre-LayerNorm normalization to the queries, keys, and values, and does not include any positional embeddings in this cross-attention component. A residual connection is also added around the cross-attention layer for further stability.

\subsection{Local Decoder} 

The local decoder $\mathcal{D}$, like its encoder counterpart, utilizes a lightweight transformer architecture composed of $l_{\mathcal{D}}$ layers, where $l_{\mathcal{D}} << l_{\mathcal{G}}$. Its role is to reconstruct the original sequence of raw bytes $y_i$ from the global patch representations $o_j$. This reconstruction is carried out by modeling the conditional dependencies between bytes, utilizing as input the hidden states generated by the local encoder. The architecture alternates $l_{\mathcal{D}}$ times between cross-attention modules and transformer layers. In each iteration, the cross-attention mechanism precedes the transformer layer, allowing the decoder to derive byte-wise representations from the higher-level patch representations, which are then processed further by the subsequent local transformer layer operating over the byte sequence.

\subsubsection{Decoder Multi-headed Cross-Attention} 

In the decoder cross-attention, the roles of the queries and key/values are interchanged i.e. the byte-representations are now the queries, and the patch representations are now the key/values. The initial byte-representations for the cross-attention are initialized as the byte embeddings from the last encoder layer i.e. $h_{l_{\mathcal{E}}}$. The subsequent byte-representations for layer $l$, $d_{l,i}$ are computed as:

\begin{align}
D_0 &= h_{l_{\mathcal{E}}} \\
B_l &= D_{l-1} + W_o\left(\text{softmax}\left(\frac{QK^T}{\sqrt{d_k}}\right)V\right), \\
  &\text{where } Q_i = W_q(d_{l-1,i}), K_i = W_k(\mathcal{D}_C(o_j)), V_i = W_v(\mathcal{D}_C(o_j)) \\
  D_l &= \text{Decoder-Transformer-layer}_l(B_l)
\end{align}

In this context, $W_k$ and $W_v$ denote the key and value projection matrices, which perform linear transformations followed by the splitting operation $\mathcal{D}_C$ on the terminal patch encodings $o_j$ generated by the global model. The matrix $W_q$ serves as the query projection and is applied to the byte-level representations $d_{l-1}$ output by the preceding decoder transformer layer (or to $h_{l_{\mathcal{E}}}$ when operating in the initial layer). Meanwhile, $W_o$ acts as the final output projection matrix, resulting in $B$ having shape $\mathbb{R}^{h_{\mathcal{D}} \times n_b}$, where $n_b$ indicates the count of output bytes. To obtain the subsequent decoder representations $D_l$, a decoder transformer layer processes the output $B$ from the cross-attention mechanism. Consistent with the procedure for local encoder cross-attention, this setup employs multi-headed attention, employs pre-LayerNorm configurations, omits positional embeddings, and utilizes a residual connection encapsulating the cross-attention module.

\section{Experimental Setup}
\label{section:experiments}
To rigorously evaluate {\textsc BLT} relative to models based on tokenization, we construct tightly controlled experiments, ensuring that {\textsc BLT} is not afforded any undue benefits from handling potentially longer sequence inputs.

\subsection{Pre-training Datasets} In our experiments across different model scales, we utilize two primary datasets for pre-training: (1) The Llama 2 dataset~\citep{touvron2023llama}, which consists of 2 trillion tokens sourced from a wide range of publicly accessible materials, subsequently subjected to rigorous cleaning and filtering procedures to ensure higher data quality; and (2) {\textsc BLT}-1T, a newly constructed corpus comprising 1 trillion tokens from diverse public domains, incorporating, as well, a subset of the Datacomp-LM pre-training dataset~\citep{li2024datacomp}. The Llama 2 dataset underpins scaling law analyses, specifically concerning the optimal token count identified in~\cite{dubey2024llama}, guiding architectural optimization for {\textsc BLT}. In contrast, the {\textsc BLT}-1T dataset is employed in comprehensive pre-training in order to evaluate downstream performance relative to Llama 3. Importantly, neither dataset contains any information derived from Meta’s products or services. For experiments involving tokenizer-based models, the Llama 3 tokenizer—featuring a vocabulary of 128K tokens—is adopted, as it consistently yielded superior baseline results compared to the Llama 2 tokenizer in our study.

\subsection{Entropy Model} The entropy model utilized throughout our experiments consists of a byte-level language model that is trained on the identical training set as the full {\textsc BLT} model. By default, the architecture adopted is a transformer comprising 100M parameters, organized in 14 layers, each with a hidden size of 512, and featuring a sliding window attention mechanism spanning 512 bytes. All other hyperparameter configurations align with those applied to our local and global transformer variants. Additional experiments explored variations in model scale, receptive field, and architecture, as outlined in \autoref{section:ablations}. Notably, for sufficiently limited receptive fields, the learned entropy model lends itself to representation as a compact lookup table.

\subsection{Entropy Threshold and Equalizing Context Length} For approaches relying on entropy-based patching, we determine a patching threshold that results in the intended mean \textit{patch size} across the pretraining data mixture. Unlike conventional tokenization, {\textsc BLT} permits arbitrary selection of \textit{patch size}, which can considerably impact the effective context window accessible to the model. To ensure fairness and prevent models with larger patch sizes from accessing greater context, we fix the expected number of bytes per batch, keeping it unchanged across different configurations. As a result, when employing larger patch sizes, the model’s sequence length is proportionally decreased. For Llama 2 data, we use an 8k byte context window, whereas with the {\textsc BLT}-1T dataset, we set the average context window to 16k bytes, holding the overall average batch size at 16M bytes.

Although the average batch size remains unchanged, dynamic patching techniques produce varying proportions of bytes per patch during the data loading process. To optimize computational efficiency, our approach to {\textsc BLT} training involves grouping batches based on the number of patches, thereby eliminating the need for padding within the latent transformer, which is comparatively resource-intensive. This approach guarantees a consistent number of patches in every batch. To further manage memory usage and prevent surges caused by exceptionally large patches, byte sequences are padded or, when necessary, truncated to lengths of 12k bytes for Llama 2 and 24k bytes for the {\textsc BLT}-1T datasets throughout training.

\subsection{Entropy Model Context}
\label{section:entropy-context}
Our observations indicate that entropy patching tends to generate increasingly larger patches when applied to structured datasets such as multiple choice tasks, which are characterized by high repetitiveness (refer to the MMLU example patching in \autoref{fig:mmlu_patching}). These differences arise due to the entropy model context assigning lower entropy values to repeated content. Consequently, for the large-scale experiment with {\textsc BLT}-Entropy and a patch size of 4.5, we refresh the entropy context by inserting new lines and employ an approximate monotonicity constraint, as this approach is less prone to issues of "entropy drift" stemming from variations in context length. This adjustment only influences our entropy computation methodology; our approach to determining the entropy threshold remains unchanged.

\subsection{FLOPs Estimation} 
\label{section:flops}

Our approach to calculating transformer floating point operations is primarily based on the methodology described in Chinchilla~\citep{hoffmann2022training}. This includes tallying the flops attributed to feed-forward networks, the qkvo projections within the self-attention module, as well as the attention computation and the output projection. In contrast to the referenced method, we make the assumption that the input embedding layer leverages an efficient lookup implementation rather than employing dense matrix multiplications, which renders this step as incurring zero flops. Consistent with established convention, we approximate that the backward pass entails twice the number of flops when compared to the forward pass.

To compute flops \textit{per byte} for {\textsc BLT} models, we add up the flops for the local encoder transformer, the global latent transformer, and the local decoder transformer, together with the cross attention blocks in the encoder and the decoder:

\begin{align}
    \text{FL}_{\text{{\textsc BLT}}} &= \frac{1}{n_p} \, \text{Transf. FL}(h_{\mathcal{G}}, l_{\mathcal{G}}, m=n_{ctx}/n_p, V=0) \\
   &+ \text{Transf. FL}(h_{\mathcal{E}}, l_{\mathcal{E}}, m=w_{\mathcal{E}}, V=0) \\
   &+ \text{Transf. FL}(h_{\mathcal{D}}, l_{\mathcal{D}}, m=w_{\mathcal{D}}, V=256) \\ 
    &+ \frac{k}{n_p} \, \text{Cross Attn. FL}(h_{\mathcal{E}}, l_{\mathcal{E}}, m=n_p, r=n_p/k) \\ 
   &+ \text{Cross Attn. FL}(h_{\mathcal{D}}, l_{\mathcal{D}}, m=k, r=k/n_p) 
\end{align}

Here, $n_{ctx}$ denotes the number of bytes in the input sequence, $n_p$ specifies the patch size, $r$ represents the proportion of queries relative to keys and values, and $k$ describes the relationship between patch and byte dimensions—namely, the count of local model partitions that are concatenated to yield the complete global representation, as illustrated with $k=2$ in \autoref{fig:crossattn}. The parameter $V$ indicates the size of the vocabulary for the output projection, relevant exclusively for the local decoder component. Depending on whether a particular module operates on the byte sequence or on the patch sequence, the attention mechanism adapts its context length, denoted $m$. We adjust the attention computation flops for each architectural module accordingly. Detailed formulations for calculating flops in both the Transformer and Cross-Attention components can be found in Appendix~\ref{section:flops-appendix}.

\subsection{Bits-Per-Byte Estimation} Perplexity only makes sense in the context of a fixed tokenizer as it is a measure of the uncertainty for each token. When comparing byte and token-level models, following previous work~\citep{xue2022byt5, yu2023megabyte, wang2024mambabyte}, we instead report Bits-Per-Byte (BPB), a tokenizer independent version of perplexity. Specifically:

In this formulation, the aggregate cross-entropy loss across the data $\pmb{x}$ quantifies the overall uncertainty, which is then scaled by dividing by both the number of bytes present in $\pmb{x}$ and the constant factor $\ln(2)$.

\subsection{Transformer Architecture Hyperparameters} Across all transformer components of {\textsc BLT}, including both the local and global modules, we adopt architectural choices consistent with those used in Llama~3~\citep{dubey2024llama}. Specifically, our feed-forward blocks incorporate the SwiGLU activation function~\citep{swiglu}, while rotary positional embeddings (RoPE)~\citep{RoPe} with a base $\theta=500000$~\citep{xiong2024effective} are applied solely within the self-attention mechanisms. For layer normalization, we leverage RMSNorm~\citep{RMSNorm}. Fixed-standard attention masking schemes, such as \textit{block causal} and \textit{fixed-window block causal}, are implemented with Flash attention~\citep{dao2022flashattention}, where the attention window size is set to 512 for masks with fixed width. When processing cross-attention that relies on masks varying adaptively based on input patches, we employ Flex Attention\footnote{\url{https://pytorch.org/blog/flexattention}}, which provides efficient, fused implementations and substantially accelerates training.

\subsection{{\textsc BLT}-Specific Hyperparameters} To evaluate the capabilities of the {\textsc BLT} models, we design experiments focusing on two main axes: scaling behavior and downstream application performance, considering model sizes of 400M, 1B, 2B, 4B, and 8B parameters. Details on the specific architectural hyperparameters corresponding to each model are provided in Appendix~Table~\ref{tab:arch}. 
For initializing queries in the first cross-attention layer of the local encoder, max-pooling is employed. Each {\textsc BLT} model is configured with $500,000$ hashes, utilizing a single hash function, and supports n-gram ranges from 3 to 8. A consistent learning rate of $4e-4$ is applied to all models in the study. This learning rate was selected by conducting a hyperparameter sweep between $1e-3$ and $1e-4$ at the 400M and 1B scales, which indicated that this value was optimal for both token-based and {\textsc BLT} variants. Training batch sizes for scaling experiments on Llama-2 datasets are aligned with the recommendations of~\cite{dubey2024llama}, or are matched by byte-size as appropriate. Model optimization uses the AdamW algorithm~\citep{adamw}, setting $\beta_1$ at 0.9, $\beta_2$ at 0.95, and $\epsilon=10^{-8}$. A learning rate schedule with a linear warm-up spanning 2000 steps is followed by a cosine decay to zero, with a weight decay factor set at 0.1, and global gradient norms clipped to a maximum of 1.0.

\section{Scaling Trends}
\label{section:scaling}

We offer a comprehensive overview of the scaling behaviors observed in byte-level models, which can provide valuable guidance for the continued development of {\textsc BLT} models. Our scaling analysis addresses the shortcomings found in earlier studies of byte-level approaches in several key respects: (a) We analyze performance under compute-optimal training conditions, (b) We train equivalent 8B-parameter models using substantial datasets (reaching up to 1T tokens or 4T bytes) and assess their capabilities on downstream evaluation tasks, and (c) We quantify scaling dynamics under settings where inference costs are kept constant. In subsequent sections, we will examine particular benefits that arise from modeling sequences at the byte level.

\subsection{Parameter Matched Compute Optimal Scaling Trends}
\label{section:parameter-matched-scaling}
We conduct experiments using the Llama 2 dataset to train several \textit{compute-optimal} bpe and {\textsc BLT} models spanning four scales, from 1B up to 8B parameters. Training compute, measured in flops, is then compared to language modeling accuracy on a representative sample of the data mixture. The bpe models are trained following the parameter-to-data ratio optimized by Llama~3~\citep{dubey2024llama}, which identifies the setting that delivers optimal performance for a fixed compute budget~\citep{hoffmann2022training}. This configuration serves as a rigorous reference for assessing our models. For every bpe model, a corresponding {\textsc BLT} model is also trained on identical data, employing a Latent Transformer with matching model size and architectural details to ensure a fair comparison.

As shown in~\autoref{fig:scaling-compute-opt} (right), {\textsc BLT} models consistently achieve performance comparable to or better than their bpe equivalents, and this pattern remains consistent across increasing model sizes and computational budgets. To our knowledge, {\textsc BLT} represents the first byte-level Transformer architecture to reach scaling behavior on par with BPE-based models under compute-optimal conditions. This finding substantiates our hypothesis that the ideal balance between parameter count and training compute for bpe extends to {\textsc BLT}, or deviates only minimally.

Achieving competitive bpe scaling patterns necessitates both advances in architecture and the adoption of dynamic patching techniques. In ~\autoref{fig:scaling-compute-opt} (left), we present a comparison between models utilizing space-patching and Llama 3. To approximate the performance of SpaceByte \citep{slagle2024spacebyte}, we employ the {\textsc BLT} space-patching approach, omitting n-gram embeddings and cross-attention mechanisms. While SpaceByte demonstrates enhancements over Megabyte, its performance still falls short relative to Llama 3. The right panel of \autoref{fig:scaling-compute-opt} displays the performance gains attributable to the integration of architectural innovations and dynamic patching. At large scales, {\textsc BLT} models achieve results comparable to the leading tokenizer-based architectures such as Llama 3.

Additionally, we note that the selection of tokenizer impacts model performance in tokenizer-dependent architectures; specifically, models utilizing the Llama-3 tokenizer achieve superior results compared to those employing the Llama-2 tokenizer, even when both are trained on identical datasets.

In summary, the {\textsc BLT} model exhibits scaling dynamics that fall between those of Llama 2 and Llama 3 when patch sizes are substantially increased. Specifically, Llama 2 and Llama 3’s bpe tokenizers yield average token lengths of 3.7 and 4.4 bytes, respectively, whereas {\textsc BLT} demonstrates comparable scaling behavior with average patch dimensions of 6 or even 8 bytes. Since the number of inference flops required decreases as the average patch size grows, adopting an 8-byte patch size translates to almost a 50\% reduction in inference flops. Empirically, models utilizing these larger patch sizes also display enhanced performance with increasing model and dataset scale. Although {\textsc BLT} initialized with 8-byte patches lags behind the bpe-encoded Llama 2 at the 1B parameter mark, it surpasses bpe when scaled to 7B parameters. These results indicate that as models and training budgets expand, patch sizes of this magnitude—or potentially even greater—could lead to further improvements in performance.

\subsection{Beyond Compute Optimal Task Evaluations}
\label{section:evals}
To investigate scaling dynamics more deeply, we train an 8B {\textsc BLT} model past the compute-optimal regime utilizing the {\textsc BLT}-1T dataset, which is both larger and of higher quality. We then evaluate its capabilities on a diverse array of established classification and generation benchmarks. For this assessment, we choose tasks aimed at measuring common sense reasoning, global knowledge, and code synthesis:
\paragraph{Classification tasks} comprise ARC-Easy (0-shot)~\citep{Eval_ARC}, Arc-Challenge (0-shot)~\citep{Eval_ARC}, HellaSwag (0-shot)~\citep{Eval_hellaswag}, PIQA (0-shot)~\citep{Eval_piqa}, and MMLU (5-shot)~\citep{hendrycks2020measuring}. Our evaluation procedure applies a prompt-scoring technique based on computing likelihoods across multiple-choice responses, and we aggregate the mean accuracy over the selected tasks.
\paragraph{Coding related generation tasks:} We evaluate the model’s code generation proficiency by reporting pass@1 scores on both MBPP (3-shot)~\citep{Eval_MBPP} and HumanEval (0-shot)~\citep{Eval_HumanEval}, measuring the model’s success rate at generating correct Python code on the first attempt.

\autoref{tab:evals} presents a comparison among three models trained on the {\textsc BLT}-1T dataset: a model that utilizes the Llama 3 tokenizer with bpe encoding,\footnote{The Llama 3 tokenizer, featuring a 128k vocabulary, is selected due to its superior performance relative to Llama 2's 32k vocabulary.} alongside two {\textsc BLT} model configurations. The first variant incorporates a space-patching mechanism ({\textsc BLT}-Space), while the second leverages an entropy-based patching strategy ({\textsc BLT}-Entropy), both subject to an approximate monotonicity constraint and with the context for the entropy-based approach being reset at new line characters (elaborated in~\autoref{section:entropy-context}). Training of all three models is carried out under an equivalent computational (flop) budget. Notably, for {\textsc BLT}-Entropy, we further lower the entropy threshold at inference from 0.6 to 0.1—a modification that, although increasing the number of inference steps, leads to enhanced task performance.

The {\textsc BLT}-Entropy model achieves superior performance compared to the Llama 3 model on four out of seven tasks, despite both models being trained with an equivalent volume of data. This advantage likely stems from two main factors: (1) more efficient utilization of training computation enabled by dynamic patching, and (2) the explicit modeling of data at the byte level rather than at the token level.

Conversely, {\textsc BLT}-Space demonstrates lower performance relative to the Llama 3 tokenizer in all tasks except one. Nevertheless, its average patch size of 6 bytes yields a considerable decrease in inference flops. By contrast, the models utilizing bpe and entropy-patching approaches display an average patch size near 4.5 bytes on the training dataset. Given an identical training budget, the larger patch size allows the model to process 30\% more data compared to the other two, a factor that may cause {\textsc BLT} to deviate further from compute-optimal efficiency.

\subsection{Patches Scale Better Than Tokens} In {\textsc BLT} models, it is possible to enlarge both the model itself and the patch size at the same time, all without exceeding a fixed computation budget for training and inference or altering the amount of training data. This flexible increase in patch size sets patch-based approaches apart from conventional token-based models, which are constrained by the efficiency compromises of a static vocabulary as outlined in Section~\ref{section:bpe-patch}. Using longer patches reduces computational demands, and the saved resources can be directed toward expanding the global latent transformer, as it needs to be executed less frequently.

We perform a scaling experiment at constant inference cost to investigate whether increasing model size, coupled with fewer inference steps on larger input patches, leads to improved performance relative to smaller models operating on smaller patches with more steps. Using base models of 400m and 3.6B parameters alongside the Llama 2 tokenizer, we identify models with comparable computational budgets (in FLOPs) for both the Llama 3 tokenizer and {\textsc BLT}-Entropy models, the latter using average patch sizes of 6 and 8 bytes on the training data mixture (refer to~\autoref{tab:fixed-inf-params} for specific model configurations). Notably, models utilizing patch size 8 are equipped with 3 encoder layers rather than 1. Each configuration is trained across different training FLOP allocations.

As illustrated in \autoref{fig:fixed_inference_scaling}, {\textsc BLT} models exhibit superior scaling performance compared to tokenization-based approaches across both inference flop categories. For lower training resource scenarios, BPE models initially outperform, but their advantage is soon overtaken by {\textsc BLT} models once surpassing the compute-optimal threshold. In practical terms, allocating increased computational resources during the initial pretraining phase can yield more capable models under a fixed inference computational budget. An exemplary case is the 8B model category, such as Llama 3.1, which has undergone training with data volumes approximately two orders of magnitude greater than the compute-optimal amount for that particular model size.

The crossover point at which {\textsc BLT} surpasses token-based models has moved slightly nearer to the compute-optimal boundary in the case of larger flop class models, shifting from about three times to 2.5 times the compute-optimal budget. Likewise, the model utilizing a patch size of 8 demonstrates a sharper scaling curve within the higher flop regime, enabling it to exceed the performance of other models at an earlier stage. As noted in Section~\ref{section:parameter-matched-scaling}, models with increased patch sizes tend to approach the performance of BPE models as model scale increases. We partially attribute this phenomenon to the fact that the proportion of overall flops consumed by the byte-level Encoder and Decoder diminishes at larger scales, given their slower growth relative to the Latent Transformer. When scaling the total number of parameters by a factor of 20 (from 400M to 8B), the number of local model parameters within {\textsc BLT} increases only by approximately twofold. This observation is relevant because increases in patch size impact solely the flops in the patch Latent Transformer, leaving those in the byte-level modules unaffected. Consequently, {\textsc BLT}-Entropy with patch size 8 increased from 1.6x to 1.7x the Llama 2 model size when progressing to the larger scale model.

In conclusion, our investigation into patch-length scaling reveals that the {\textsc BLT} patch-based framework exhibits superior scaling performance when both the patch size and the model size are increased together. These scaling behaviors appear to hold and may even be enhanced as the model is scaled up further.

\section{Byte Modeling Improves Robustness} Furthermore, we evaluate the robustness of {\textsc BLT} relative to token-based architectures that do not natively incorporate byte-level details, and introduce a method to convert pretrained token-based models into their byte-ified counterparts.
\label{sec:robustness}

\subsection{Character-Level Tasks}

Initial interest in training models at the byte level arose from their heightened resistance to input noise at the byte scale and their intrinsic understanding of token composition—areas in which conventional tokenizer-based approaches are less effective. To empirically assess these qualities, we conduct supplementary evaluations on tasks specifically designed to test robustness to input corruption as well as sensitivity to character information. These assessments encompass English and multi-lingual characters, as well as digits and phonemes. Table~\ref{tab:char_tasks} summarizes the findings from these character-focused tasks.

\paragraph{Noisy Data} To assess the resilience of tokenizer-based models in contrast to {\textsc BLT}, we generate altered versions of the benchmark classification datasets outlined in Section~\ref{section:evals}. Specifically, we implement five separate character-level perturbation techniques to modify the text: (a) \textit{AntSpeak}: This method reformats the text entirely into uppercase characters, each separated by spaces. (b) \textit{Drop}: Randomly eliminates 10\% of the characters from the text. (c) \textit{RandomCase}: Alters the case of the characters by randomly assigning 50\% to uppercase and 50\% to lowercase. (d) \textit{Repeat}: Selects 20\% of characters and repeats them as many as four times. (e) \textit{UpperCase}: Converts every character in the text to uppercase. For evaluation, each noising technique is applied individually to either the prompt, the completion, or both components, treating each configuration as a separate task; the results are averaged. Table \ref{tab:char_tasks} presents the performance on noised HellaSwag~\citep{Eval_hellaswag}, showing that {\textsc BLT} consistently surpasses tokenizer-based models in robustness, with an average improvement of 8 points relative to a similarly trained model, and even exceeds the performance of the Llama 3.1 model, which was trained with substantially more data.

\paragraph{Phonology - Grapheme-to-Phoneme (G2P)} We evaluate {\textsc BLT}'s proficiency in converting a string of graphemes—written characters forming a word—into its corresponding phonemic representation, which reflects the word's pronunciation. Table \ref{tab:char_tasks} shows the G2P task outcomes under a 5-shot regime, utilizing the Phonology Bench~\citep{suvarna2024phonologybench}. Our results indicate that {\textsc BLT} achieves superior performance compared to the Llama 3 1T tokenizer-based baseline on this task.

\paragraph{CUTE}
To evaluate proficiency at the character level, we apply {\textsc BLT} to the CUTE benchmark~\citep{edman2024cute}, which involves a diverse suite of tasks broadly sorted into three types: composition comprehension, orthographic similarity, and sequence manipulation. This suite is notably difficult for most models utilizing tokenization, as these models tend to recognize the orthographic properties of their tokens but lack the flexibility to deploy such knowledge for manipulating text. As presented in Table~\ref{tab:char_tasks}, {\textsc BLT}-Entropy significantly outperforms both Llama 3 BPE models on CUTE, exceeding their performance by over 25 points. Notably, our approach excels at character manipulation, achieving 99.9\% accuracy on the spelling-focused tasks. These substantial gains, despite the fact that {\textsc BLT} was trained on just one-sixteenth of the data used for Llama 3.1, underscore the challenge BPE architectures face in acquiring fine-grained character information. Figure \ref{fig:cute_examples} showcases instances in which Llama 3’s tokenizer-based model falters, whereas our {\textsc BLT} exhibits strong performance. Word deletion and insertion are the only categories where BPE-based models hold an advantage, likely due to the inherent difficulty that byte-level models encounter with such operations; however, the margin is relatively narrow, suggesting that constructing word-level manipulations from character representations could be more straightforward than disassembling token-level structure. Evaluation protocols were standardized across all tasks, using original Huggingface prompts. It is worth noting that BPE-based models could potentially improve with further prompt optimization.

\paragraph{Low Resource Machine Translation} We assess the performance of {\textsc BLT} on translation tasks involving both directions across six major language families and twenty-one low-resource languages, each utilizing distinct scripts, as defined by the FLORES-101 benchmark~\citep{goyal2022flores}. The SentencePiece BLEU results are presented in Table~\ref{tab:flores}. {\textsc BLT} exhibits superior results compared to a system trained with the Llama 3 tokenizer, providing an overall improvement of 2 BLEU points for translations into English and 0.5 points for translations from English. For well-studied language pairs, {\textsc BLT}'s performance is on par with or marginally surpasses that of Llama 3. More notably, in the context of lower-resource language families, {\textsc BLT} consistently outperforms Llama 3, highlighting byte modeling's capability to better generalize to rare and diverse byte sequences.

\subsection{Training {\textsc BLT} from Llama 3}
\label{sec:distillation}
We investigate an approach where {\textsc BLT} models benefit from the initialization of their global transformer parameters with those from a pre-trained Llama 3.1 model, thereby utilizing well-established tokenizer-based models to accelerate and improve training convergence. After this initialization, we proceed by updating the global transformer’s weights using a learning rate that is one-tenth of that used for the local encoder and decoder components, following the same number of optimization steps as in Llama 3. The resulting performance, compared against a baseline {\textsc BLT}, is reported in Table~\ref{tab:distill}. The findings indicate that {\textsc BLT} models initialized from Llama 3.1 parameters markedly surpass both the Llama 3 and standard {\textsc BLT} baselines, with all models trained under equivalent computational budgets. Additionally, relative to the {\textsc BLT}-Entropy variant in Table~\ref{tab:evals}—which underwent training on a much larger corpus of 1T tokens—{\textsc BLT} initialized from Llama 3.1 demonstrates improved results on the MMLU benchmark. These observations highlight the potential of this initialization strategy for substantially decreasing training flops while maintaining or improving performance.

This approach can alternatively be understood as adapting tokenizer-based architectures into tokenizer-free formats, thereby enabling a pre-trained LLaMA 3.1 model to function as a {\textsc BLT} model. For a thorough evaluation, Table \ref{tab:distill} presents the performance of the original LLaMA 3.1, which was pre-trained on 15T tokens, alongside results from the {\textsc BLT} version distilled from LLaMA 3. We observe that our model shows a modest decrease in accuracy on MMLU and HumanEval benchmarks, with more pronounced losses evident in other assessment areas. These observations highlight the necessity for future research aimed at more effectively harnessing the capabilities of the pre-trained model, especially through better tuning of data compositions and other key hyperparameters.

\section{Ablations and Discussion}
\label{section:ablations}
Here, we present ablation studies to motivate specific architectural decisions in {\textsc BLT} as well as to examine the patching strategy and hyper-parameter selections utilized for the {\textsc BLT} model with 8B parameters, trained on the {\textsc BLT}-1T dataset.

\paragraph{Entropy Model Hyper-parameters} In order to analyze how entropy model capacity and context window size influence scaling behavior, we conduct experiments with byte-level entropy transformer models that vary in size from 1 million to 100 million parameters and use context window lengths ranging from 64 up to 512 tokens. The results, shown as bpb versus training flop scaling law curves for our $400m$ and $1b$ {\textsc BLT} models trained on the Llama-2 dataset, are depicted in \autoref{fig:entropy_ablation}. The data indicate a positive association between both the number of parameters and context window size of the entropy model and scaling performance, though gains tend to plateau as the model size surpasses 50 million parameters.

\paragraph{Types of Patching}

We conduct an ablation study on the four patching strategies described in Section~\ref{section:patching}: (1) Strided Patching with strides of 4 and 6, (2) Patching at whitespace positions, (3) BPE Tokenizer patching utilizing the Llama 3 tokenizer, and (4) Entropy-based patching implemented with a compact byte-level language model.

Although dynamic patching shortens the actual sequence lengths, we ensure that the context length remains consistent across all patching strategies by regulating sequence length. This approach guarantees that, on average, every model processes the same quantity of bytes per sequence in both training and inference phases, thereby eliminating biases related to differing context sizes. Figure~\ref{fig:scaling-compute-opt} presents the findings from these ablation studies. Notably, every patching scheme except for static patching achieves superior performance, with space patching offering results nearly matching those of dynamic entropy-based patching.

In \autoref{tab:evals_ablation}, we report the results of benchmark tests for {\textsc BLT} models, contrasting approaches based on tokenizers, space patching, and entropy-driven patching, all trained on the Llama 2 dataset for a number of steps chosen to optimize performance~\citep{dubey2024llama}. While space patching is comparatively straightforward—eschewing the use of a real-time entropy model during the training process—our findings indicate that the improvements observed with entropy-based patching in scaling analyses (see Section~\ref{section:scaling}) are also evident on downstream benchmark evaluations.\footnote{Space patching outcomes are derived from earlier experiments conducted without cross-attention; however, similar tendencies emerge even when cross-attention is included.}

\paragraph{Cross-Attention} 
In~\autoref{tab:crossattn_ablations_1b}, we analyze the impact of adding cross-attention at different stages within the encoder and decoder components of {\textsc BLT}. For encoder cross-attention, we evaluate three query initialization methods: 1) a shared learned embedding for all global states, 2) a hash-based embedding derived from the bytes of the patch, and 3) using a pooled version of the encoder’s hidden states corresponding to the patch bytes at the relevant encoder layer.

Our results indicate that incorporating cross-attention within the \textit{decoder} yields the best performance. In the encoder, cross-attention leads to modest gains, though these benefits arise only when queries are initialized through pooling. Furthermore, our analysis reveals that cross-attention proves especially beneficial for Common-Crawl data, with its impact being more pronounced as patch sizes increase.

\paragraph{n-gram Hash Embeddings} 

We experiment with n-gram hash embedding vocabularies set at 0, 100K, 200K, and 400K, with the corresponding outcomes detailed in Table~\ref{tab:ngram_ablations_1b}. Our results demonstrate that hash embeddings consistently yield benefits across all tested domains, with the most pronounced improvements on Wikipedia and Github (where we observe a 0.04 bpb gain, versus just a 0.01 bpb difference after 15k training steps at the 8B model scale). At the 8B parameter scale, reducing hashes from 500K to 300K results in only a marginal difference—about 0.001 bpb—after 15k steps. These findings suggest that the introduction of hash embeddings is essential for elevating the performance of {\textsc BLT} to be comparable with tokenizer-based architectures, but beyond 300K hashes, additional increases produce limited benefits. Furthermore, our observations indicate that these performance gains largely supplement those achieved by cross-attention, as both methods enhance results on distinct datasets.

\paragraph{Local Model Hyperparamaters} Table \ref{tab:local_layers} presents an ablation study exploring different configurations for the number of layers in the local encoder and decoder. Results indicate that when hash n-gram embeddings are utilized, {\textsc BLT} performs effectively even with a minimal encoder—consisting of only a single layer—while benefiting from a more complex, deeper decoder architecture.

\section{Related Work}
\label{section:related_work}

\textbf{Character-Level RNNs:} Since the early development of neural language models~\citep{sutskever2011generating, mikolov2012subword, graves2013generating}, Character Language Modeling has maintained its appeal due to its inherent capability to handle out-of-vocabulary words naturally, eliminating the need for traditional back-off techniques. \cite{kim2016character} present an architecture that exclusively processes input at the character level via convolutional and highway networks, which in turn provide inputs to LSTM-based RNNs. This approach achieves parity with the then-leading RNN-based language models on English, and surpasses them in morphologically complex languages—an important benefit of character-level large language models. Employing a similar philosophy, \cite{kenter2018byte} leverage byte-level LSTM models for machine comprehension, achieving superior results compared to word-level systems particularly for Turkish and Russian, both of which are morphologically rich. In a related vein, \cite{zhang2015character} propose character-based convolutional models for classification tasks, demonstrating that they outperform word-level models on specific benchmarks. The work by \citet{chung2019hierarchical} further enhances character-level language modeling by introducing hierarchical LSTM models incorporating boundary detectors at each stage to uncover the underlying text hierarchy. Departing from attention mechanisms, ByteNet~\citep{kalchbrenner2016neural} applies convolutional layers over characters for machine translation tasks.

\textbf{Character-Level Transformers:} The introduction of transformer architectures leveraging attention mechanisms~\citep{vaswani2017attention} and enhanced by subword tokenization strategies~\citep{sennrich-etal-2016-neural} marked a substantial advancement in the capabilities of neural language models for both language modeling and various benchmark evaluations. Nevertheless, segmenting text into words or subwords inherently imposes a specific inductive bias regarding the granularity at which models process linguistic information. Seeking to integrate the strengths of transformer models with encouraging findings from character-level language modeling, \citet{al2019character} employed very deep transformer networks, augmented with auxiliary loss functions. These approaches enabled the training of transformer-based models that surpassed the performance of earlier character-level LSTMs. Despite these improvements, the performance gap between character-level and word-level LLMs remained notable. GPT-2~\citep{radfordgpt2} further demonstrated that on large-scale corpora, such as the 1 Billion Word Benchmark, byte-level language models still lagged behind those based on word units.

\cite{Choe2019BridgingTG} found that transformer-based large language models (LLMs) operating at the byte level can surpass subword-level LLMs with equivalent parameter counts in terms of performance; however, these byte-level models incur substantially higher computational costs and require longer training times. In a similar vein, \cite{el2020characterbert} introduced CharFormer, a variant of BERT that derives word representations via convolutional layers on character embeddings, yielding performance gains in the biomedical field, albeit at the expense of much greater computational resources. \cite{clark2022canine} propose CANINE, an encoder-only architecture with 150M parameters designed to process raw character inputs. At its core lies a deep transformer component reminiscent of the global model described herein, and its architecture integrates both a local transformer and strided convolution operations to downsample character input. CANINE demonstrates superior downstream results compared to mBERT, a similarly sized token-level encoder-only baseline, on multilingual benchmarks. ByT5~\citep{xue2022byt5} examined byte-level encoder-decoder architectures, consciously avoiding patch-based reductions. Although this model displayed greater robustness to input noise and matched the performance of tokenizer-based models with a fourth of the training data, the absence of patching necessitated performing computation-intensive attention across each byte position, which greatly inflated the computational burden. Direct byte-level modeling increases input sequence lengths, which hinders efficient scaling of such models. More recently, \cite{wang2024mambabyte} leveraged the Mamba Architecture~\citep{gu2023mamba}, featuring a mechanism for maintaining a constant memory footprint even over long contexts, to construct a byte-level Mamba model—again eschewing patching—and demonstrated that, under flop-matched comparisons at 350M parameters, their approach outperformed byte-level transformer counterparts in terms of bits-per-byte metrics across several corpora.

\textbf{Patching-based approaches:} The utilization of patching techniques can significantly reduce the computational overhead (measured in flops) typically associated with byte-level LLMs, all while maintaining competitive performance. Multiple studies have demonstrated promising results on limited model sizes and data volumes. For instance, \cite{nawrot-etal-2022-hierarchical} introduce the hourglass transformer by employing static patching strategies—combining downsampling and upsampling—which shows superior results compared to other byte-level benchmarks at the 150M parameter level. Building on this, \cite{nawrot-etal-2023-efficient} incorporate dynamic patching methods, including an end-to-end trained boundary-predictor, a boundary-predictor guided by supervision from specific tokenizers, and an entropy-driven patching model analogous to {\textsc BLT}. Their findings indicate that these techniques enable models to surpass contemporaneous vanilla transformers in language modeling tasks at the 40M parameter scale trained on approximately 400M tokens. Moreover, \cite{lester2024training} explore models trained on inputs compressed via arithmetic coding, achieving higher compression than BPE, and introduce an equal-info windows technique. Their experiments demonstrate improvements over byte-level baselines on language modeling benchmarks, although they still lag behind subword-based approaches.

Our research takes particular inspiration from MegaByte~\citep{yu2023megabyte}, a decoder-only, causally-masked language model that implements a predetermined static patching and concatenates encoded representations to transform byte sequences into patches. This approach also employs a localized decoder module to reconstruct byte sequences from these patches. The original MegaByte work demonstrates comparable performance to tokenizer-driven models at a parameter scale of 1B on datasets comprising approximately 400B bytes. Throughout our experimental analyses, we systematically examine MegaByte and observe that static patching falls short of matching compute-efficient, state-of-the-art tokenizer-based models in a flop-constrained environment. In response, we illustrate how {\textsc BLT} addresses and narrows this performance discrepancy. \cite{slagle2024spacebyte} corroborate our findings, highlighting MegaByte’s limitations and advocating for an expanded patching strategy—specifically, segmenting based on whitespace and space-like bytes and incorporating a localized encoder. Their results indicate that this revised design can surpass tokenized transformer models on selected corpora such as Github and arXiv in the 1B parameter setting under compute controls. We include an evaluation of this modified architecture in our work, further establishing that more sophisticated architectural changes are essential for byte-level models to reach parity with leading-edge token-based models like Llama 3, especially as scale increases.

\section{Limitations and Future Work}

In this study, our approach to architectural configuration relies on training models for the number of steps found to be optimal for Llama~3~\citep{dubey2024llama}. It is important to note, however, that these scaling laws were developed specifically for BPE-level transformers, which may not produce ideal data-to-parameter ratios when applied to {\textsc BLT}. Developing and verifying scaling laws tailored to {\textsc BLT} is left as a direction for future investigation and could reveal even more advantageous scaling behavior for this architecture. Furthermore, the majority of our experiments were limited to model sizes up to 1B parameters. As we extend our work to models with 8B parameters and larger, optimal architectural configurations may shift, potentially yielding better results at these greater scales.

Current transformer codebases and libraries are optimized primarily for architectures based on tokenizers. Although we conduct experiments with matched theoretical flops and integrate select optimized components (like FlexAttention) to accommodate layers that diverge from standard transformer designs, our implementations might still lag behind tokenizer-based models regarding actual runtime and could potentially achieve better efficiency with additional optimizations.

Although {\textsc BLT} currently relies on a separately trained entropy model for its patching process, integrating the learning of this patching component within an end-to-end framework represents a promising avenue for future research. As described in Section \ref{sec:distillation}, we provide preliminary experimental results that suggest the viability of adapting tokenizer-based models like Llama 3—trained on over 10T tokens—into a "byte-ified" form by initializing the global transformer using their pre-trained weights and then freezing those parameters. Continued exploration in this line of inquiry could yield approaches that preserve the advantages associated with bytefying, and may even surpass the performance of existing tokenizer-based architectures without necessitating full retraining.

\section{Conclusion}
\label{section:conclusion}

In this work, we introduce the Byte Latent Transformer (\textbf{BLT}), an innovative model architecture that challenges the standard reliance on fixed-vocabulary tokenization in large language models. \textsc{BLT} employs an adaptive, learnable mechanism for aggregating bytes into variable-length patches, enabling the model to allocate computational effort more effectively in response to input complexity. This leads to notable gains in both processing efficiency and model resilience. Through an extensive scaling analysis, we show that \textsc{BLT} achieves performance on par with leading tokenization-based models such as Llama 3, scaling up to 8B parameters and 4T processed bytes, and is capable of exchanging slight evaluation performance drops for up to a 50\% reduction in inference computational cost. In addition, \textsc{BLT} facilitates a novel scaling dimension, allowing increases in both model capacity and patch granularity without exceeding a predetermined inference resource ceiling, making it particularly suitable for realistic computational budgets. By working directly with byte-level input, \textsc{BLT} enhances handling of rare or unconventional data and demonstrates superior robustness against input noise, along with improved representation of sub-word information. Collectively, these advances establish \textsc{BLT} as a compelling successor to traditional tokenization techniques, offering a flexible and resilient structure for building more efficient and versatile language models.

\end{document}